\chapter{Anhang: Konjugationstabellen}
\label{cha:Anhang}

Dieser Anhang enthält alle wichtigen Konjugationstabellen als Referenz.

\section{Verbkonjugation: Präsens}
\label{sec:KonjugationPraesens}

\subsection{schwache Verben (regelmäßig)}
\label{sec:SchwacheVerben}

\begin{center}
\begin{tabular}{|l|c|c|c|c|}
\hline
\textbf{} & \textbf{machen} & \textbf{sagen} & \textbf{brauchen} & \textbf{spielen} \\ \hline
ich & mache & sage & brauche & spiele \\ \hline
du & machst & sagst & brauchst & spielst \\ \hline
er/sie/es & macht & sagt & braucht & spielt \\ \hline
wir & machen & sagen & brauchen & spielen \\ \hline
ihr & macht & sagt & braucht & spielt \\ \hline
sie/Sie & machen & sagen & brauchen & spielen \\ \hline
\end{tabular}
\end{center}

\subsection{starke Verben (unregelmäßig)}
\label{sec:StarkeVerben}

\begin{center}
\begin{tabular}{|l|c|c|c|c|}
\hline
\textbf{} & \textbf{lesen} & \textbf{sprechen} & \textbf{geben} & \textbf{nehmen} \\ \hline
ich & lese & spreche & gebe & nehme \\ \hline
du & liest & sprichst & gibst & nimmst \\ \hline
er/sie/es & liest & spricht & gibt & nimmt \\ \hline
wir & lesen & sprechen & geben & nehmen \\ \hline
ihr & lest & sprecht & gebt & nehmt \\ \hline
sie/Sie & lesen & sprechen & geben & nehmen \\ \hline
\end{tabular}
\end{center}

\subsection{haben}
\label{sec:Habentaben}

\begin{center}
\begin{tabular}{|l|c|}
\hline
\textbf{} & \textbf{haben} \\ \hline
ich & habe \\ \hline
du & hast \\ \hline
er/sie/es & hat \\ \hline
wir & haben \\ \hline
ihr & habt \\ \hline
sie/Sie & haben \\ \hline
\end{tabular}
\end{center}

\subsection{sein}
\label{sec:Sein}

\begin{center}
\begin{tabular}{|l|c|}
\hline
\textbf{} & \textbf{sein} \\ \hline
ich & bin \\ \hline
du & bist \\ \hline
er/sie/es & ist \\ \hline
wir & sind \\ \hline
ihr & seid \\ \hline
sie/Sie & sind \\ \hline
\end{tabular}
\end{center}

\subsection{werden}
\label{sec:Werden}

\begin{center}
\begin{tabular}{|l|c|}
\hline
\textbf{} & \textbf{werden} \\ \hline
ich & werde \\ \hline
du & wirst \\ \hline
er/sie/es & wird \\ \hline
wir & werden \\ \hline
ihr & werdet \\ \hline
sie/Sie & werden \\ \hline
\end{tabular}
\end{center}

\section{Verbkonjugation: Präteritum}
\label{sec:KonjugationPraeteritum}

\subsection{schwache Verben}
\label{sec:SchwacheVerbenPraet}

\begin{center}
\begin{tabular}{|l|c|c|c|}
\hline
\textbf{} & \textbf{machen} & \textbf{sagen} & \textbf{brauchen} \\ \hline
ich & machte & sagte & brauchte \\ \hline
du & machtest & sagtest & brauchtest \\ \hline
er/sie/es & machte & sagte & brauchte \\ \hline
wir & machten & sagten & brauchten \\ \hline
ihr & machtet & sagtet & brauchtet \\ \hline
sie/Sie & machten & sagten & brauchten \\ \hline
\end{tabular}
\end{center}

\subsection{starke Verben}
\label{sec:StarkeVerbenPraet}

\begin{center}
\begin{tabular}{|l|c|c|c|c|}
\hline
\textbf{} & \textbf{lesen} & \textbf{sprechen} & \textbf{geben} & \textbf{nehmen} \\ \hline
ich & las & sprach & gab & nahm \\ \hline
du & lasest & sprachst & gabst & nahmst \\ \hline
er/sie/es & las & sprach & gab & nahm \\ \hline
wir & lasen & sprachen & gaben & nahmen \\ \hline
ihr & last & spracht & gabt & nahmt \\ \hline
sie/Sie & lasen & sprachen & gaben & nahmen \\ \hline
\end{tabular}
\end{center}

\section{Modalverben: Präsens}
\label{sec:ModalverbenPraesens}

\begin{center}
\begin{tabular}{|l|c|c|c|c|c|c|}
\hline
\textbf{} & \textbf{dürfen} & \textbf{können} & \textbf{müssen} & \textbf{sollen} & \textbf{wollen} & \textbf{mögen} \\ \hline
ich & darf & kann & muss & soll & will & mag \\ \hline
du & darfst & kannst & musst & sollst & willst & magst \\ \hline
er/sie/es & darf & kann & muss & soll & will & mag \\ \hline
wir & dürfen & können & müssen & sollen & wollen & mögen \\ \hline
ihr & dürft & könnt & müsst & sollt & wollt & mögt \\ \hline
sie/Sie & dürfen & können & müssen & sollen & wollen & mögen \\ \hline
\end{tabular}
\end{center}

\section{Modalverben: Präteritum}
\label{sec:ModalverbenPraeteritum}

\begin{center}
\begin{tabular}{|l|c|c|c|c|c|c|}
\hline
\textbf{} & \textbf{dürfen} & \textbf{können} & \textbf{müssen} & \textbf{sollen} & \textbf{wollen} & \textbf{mögen} \\ \hline
ich & durfte & konnte & musste & sollte & wollte & mochte \\ \hline
du & durftest & konntest & musstest & solltest & wolltest & mochtest \\ \hline
er/sie/es & durfte & konnte & musste & sollte & wollte & mochte \\ \hline
wir & durften & konnten & mussten & sollten & wollten & mochten \\ \hline
ihr & durftet & konntet & musstet & solltet & wolltet & mochtet \\ \hline
sie/Sie & durften & konnten & mussten & sollten & wollten & mochten \\ \hline
\end{tabular}
\end{center}

\section{Modalverben: Konjunktiv II}
\label{sec:ModalverbenKonjunktivII}

\begin{center}
\begin{tabular}{|l|c|c|c|c|c|c|}
\hline
\textbf{} & \textbf{dürfen} & \textbf{können} & \textbf{müssen} & \textbf{sollen} & \textbf{wollen} & \textbf{mögen} \\ \hline
ich & dürfte & könnte & müsste & sollte & wollte & möchte \\ \hline
du & dürftest & könntest & müsstest & solltest & wolltest & möchtest \\ \hline
er/sie/es & dürfte & könnte & müsste & sollte & wollte & möchte \\ \hline
wir & dürften & könnten & müssten & sollten & wollten & möchten \\ \hline
ihr & dürftet & könntet & müsstet & solltet & wolltet & möchtet \\ \hline
sie/Sie & dürften & könnten & müssten & sollten & wollten & möchten \\ \hline
\end{tabular}
\end{center}

\section{Perfekt: Hilfsverben}
\label{sec:PerfektHilfsverben}

\begin{center}
\begin{tabular}{|l|c|c|c|}
\hline
\textbf{} & \textbf{haben} & \textbf{sein} & \textbf{werden} \\ \hline
ich & habe & bin & bin \\ \hline
du & hast & bist & bist \\ \hline
er/sie/es & hat & ist & ist \\ \hline
wir & haben & sind & sind \\ \hline
ihr & habt & seid & seid \\ \hline
sie/Sie & haben & sind & sind \\ \hline
\end{tabular}
\end{center}

\begin{center}
\begin{tabular}{|l|c|c|c|}
\hline
\textbf{Partizip II} & \textbf{gemacht} & \textbf{gegangen} & \textbf{geworden} \\ \hline
\end{tabular}
\end{center}

\section{Konjunktiv I: sein}
\label{sec:KonjunktivIsein}

\begin{center}
\begin{tabular}{|l|c|}
\hline
\textbf{} & \textbf{sein (Konjunktiv I)} \\ \hline
ich & sei \\ \hline
du & seiest \\ \hline
er/sie/es & sei \\ \hline
wir & seien \\ \hline
ihr & seiet \\ \hline
sie/Sie & seien \\ \hline
\end{tabular}
\end{center}

\section{Konjunktiv I: haben}
\label{sec:KonjunktivIhaben}

\begin{center}
\begin{tabular}{|l|c|}
\hline
\textbf{} & \textbf{haben (Konjunktiv I)} \\ \hline
ich & habe \\ \hline
du & habest \\ \hline
er/sie/es & habe \\ \hline
wir & haben \\ \hline
ihr & habet \\ \hline
sie/Sie & haben \\ \hline
\end{tabular}
\end{center}

\section{Konjunktiv I: werden}
\label{sec:KonjunktivIwerden}

\begin{center}
\begin{tabular}{|l|c|}
\hline
\textbf{} & \textbf{werden (Konjunktiv I)} \\ \hline
ich & werde \\ \hline
du & werdest \\ \hline
er/sie/es & werde \\ \hline
wir & werden \\ \hline
ihr & werdet \\ \hline
sie/Sie & werden \\ \hline
\end{tabular}
\end{center}

\section{Konjunktiv II: wäre}
\label{sec:KonjunktivIIwaere}

\begin{center}
\begin{tabular}{|l|c|}
\hline
\textbf{} & \textbf{sein (Konjunktiv II)} \\ \hline
ich & wäre \\ \hline
du & wärest \\ \hline
er/sie/es & wäre \\ \hline
wir & wären \\ \hline
ihr & wärt \\ \hline
sie/Sie & wären \\ \hline
\end{tabular}
\end{center}

\section{Konjunktiv II: hätte}
\label{sec:KonjunktivIIhaette}

\begin{center}
\begin{tabular}{|l|c|}
\hline
\textbf{} & \textbf{haben (Konjunktiv II)} \\ \hline
ich & hätte \\ \hline
du & hättest \\ \hline
er/sie/es & hätte \\ \hline
wir & hätten \\ \hline
ihr & hättet \\ \hline
sie/Sie & hätten \\ \hline
\end{tabular}
\end{center}

\section{ Imperativ}
\label{sec:Imperativ}

\begin{center}
\begin{tabular}{|l|c|c|c|}
\hline
\textbf{} & \textbf{machen} & \textbf{lesen} & \textbf{sein} \\ \hline
du & mach! & lies! & sei! \\ \hline
ihr & macht! & lest! & seid! \\ \hline
Sie & machen Sie! & lesen Sie! & seien Sie! \\ \hline
\end{tabular}
\end{center}

\section{Trennbare Verben: Präsens}
\label{sec:TrennbarPraesens}

\begin{center}
\begin{tabular}{|l|c|c|c|c|}
\hline
\textbf{} & \textbf{abfahren} & \textbf{aufmachen} & \textbf{einladen} & \textbf{vorbereiten} \\ \hline
ich & fahre ab & mache auf & lade ein & bereite vor \\ \hline
du & fährst ab & machst auf & lädst ein & bereitest vor \\ \hline
er/sie/es & fährt ab & macht auf & lädt ein & bereitet vor \\ \hline
wir & fahren ab & machen auf & laden ein & bereiten vor \\ \hline
ihr & fahrt ab & macht auf & ladet ein & bereitet vor \\ \hline
sie/Sie & fahren ab & machen auf & laden ein & bereiten vor \\ \hline
\end{tabular}
\end{center}

\section{Untrennbare Verben: Präsens}
\label{sec:UntrennbarPraesens}

\begin{center}
\begin{tabular}{|l|c|c|c|c|}
\hline
\textbf{} & \textbf{erklären} & \textbf{bezahlen} & \textbf{versuchen} & \textbf{empfangen} \\ \hline
ich & erkläre & bezahle & versuche & empfange \\ \hline
du & erklärst & bezahlst & versuchst & empfängst \\ \hline
er/sie/es & erklärt & bezahlt & versucht & empfängt \\ \hline
wir & erklären & bezahlen & versuchen & empfangen \\ \hline
ihr & erklärt & bezahlt & versucht & empfängt \\ \hline
sie/Sie & erklären & bezahlen & versuchen & empfangen \\ \hline
\end{tabular}
\end{center}

\section{Wichtige unregelmäßige Verben}
\label{sec:WichtigeUnregelmaessige}

\begin{center}
\begin{tabular}{|l|c|c|c|c|}
\hline
\textbf{Infinitiv} & \textbf{Präsens} & \textbf{Präteritum} & \textbf{Perfekt} & \textbf{Bedeutung} \\ \hline
beginnen & beginnt & begann & hat begonnen & beginnen \\ \hline
bleiben & bleibt & blieb & ist geblieben & bleiben \\ \hline
bringen & bringt & brachte & hat gebracht & bringen \\ \hline
denken & denkt & dachte & hat gedacht & denken \\ \hline
fahren & fährt & fuhr & ist gefahren & fahren \\ \hline
finden & findet & fand & hat gefunden & finden \\ \hline
geben & gibt & gab & hat gegeben & geben \\ \hline
gehen & geht & ging & ist gegangen & gehen \\ \hline
gewinnen & gewinnt & gewonnt & hat gewonnen & gewinnen \\ \hline
haben & hat & hatte & hat gehabt & haben \\ \hline
halten & hält & hielt & hat gehalten & halten \\ \hline
hängen & hängt & hing & hat gehangen & hängen \\ \hline
heißen & heißt & hieß & hat geheißen & heißen \\ \hline
kennen & kennt & kannte & hat gekannt & kennen \\ \hline
kommen & kommt & kam & ist gekommen & kommen \\ \hline
können & kann & konnte & hat gekonnt & können \\ \hline
lesen & liest & las & hat gelesen & lesen \\ \hline
liegen & liegt & lag & hat gelegen & liegen \\ \hline
machen & macht & machte & hat gemacht & machen \\ \hline
mögen & mag & mochte & hat gemocht & mögen \\ \hline
müssen & muss & musste & hat gemusst & müssen \\ \hline
nehmen & nimmt & nahm & hat genommen & nehmen \\ \hline
rufen & ruft & rief & hat gerufen & rufen \\ \hline
sagen & sagt & sagte & hat gesagt & sagen \\ \hline
schlafen & schläft & schlief & hat geschlafen & schlafen \\ \hline
schlagen & schlägt & schlug & hat geschlagen & schlagen \\ \hline
schreiben & schreibt & schrieb & hat geschrieben & schreiben \\ \hline
schwimmen & schwimmt & schwamm & ist geschwommen & schwimmen \\ \hline
sehen & sieht & sah & hat gesehen & sehen \\ \hline
sein & ist & war & ist gewesen & sein \\ \hline
senden & sendet & sandte & hat gesandt & senden \\ \hline
singen & singt & sang & hat gesungen & singen \\ \hline
sitzen & sitzt & saß & hat gesessen & sitzen \\ \hline
sollen & soll & sollte & hat gesollt & sollen \\ \hline
sprechen & spricht & sprach & hat gesprochen & sprechen \\ \hline
stehen & steht & stand & hat gestanden & stehen \\ \hline
stimmen & stimmt & stimmte & hat gestimmt & stimmen \\ \hline
tragen & trägt & trug & hat getragen & tragen \\ \hline
treffen & trifft & traf & hat getroffen & treffen \\ \hline
trinken & trinkt & trank & hat getrunken & trinken \\ \hline
tun & tut & tat & hat getan & tun \\ \hline
verstehen & versteht & verstand & hat verstanden & verstehen \\ \hline
werden & wird & wurde & ist geworden & werden \\ \hline
wissen & weiß & wusste & hat gewusst & wissen \\ \hline
wollen & will & wollte & hat gewollt & wollen \\ \hline
ziehen & zieht & zog & hat gezogen & ziehen \\ \hline
\end{tabular}
\end{center}

\section{Artikel-Deklination}
\label{sec:ArtikelDeklination}

\subsection{Bestimmt}
\begin{center}
\begin{tabular}{|l|c|c|c|c|}
\hline
\textbf{Kasus} & \textbf{Maskulin} & \textbf{Feminin} & \textbf{Neutral} & \textbf{Plural} \\ \hline
Nominativ & der & die & das & die \\ \hline
Akkusativ & den & die & das & die \\ \hline
Dativ & dem & der & dem & den \\ \hline
Genitiv & des & der & des & der \\ \hline
\end{tabular}
\end{center}

\subsection{Unbestimmt}
\begin{center}
\begin{tabular}{|l|c|c|c|c|}
\hline
\textbf{Kasus} & \textbf{Maskulin} & \textbf{Feminin} & \textbf{Neutral} & \textbf{Plural} \\ \hline
Nominativ & ein & eine & ein & keine \\ \hline
Akkusativ & einen & eine & ein & keine \\ \hline
Dativ & einem & einer & einem & keinen \\ \hline
Genitiv & eines & einer & eines & keiner \\ \hline
\end{tabular}
\end{center}

\section{Adjektiv-Deklination}
\label{sec:AdjektivDeklination}

\subsection{Mit bestimmtem Artikel}
\begin{center}
\begin{tabular}{|l|c|c|c|c|}
\hline
\textbf{Kasus} & \textbf{Maskulin} & \textbf{Feminin} & \textbf{Neutral} & \textbf{Plural} \\ \hline
Nominativ & der gute & die gute & das gute & die guten \\ \hline
Akkusativ & den guten & die gute & das gute & die guten \\ \hline
Dativ & dem guten & der guten & dem guten & den guten \\ \hline
Genitiv & des guten & der guten & des guten & der guten \\ \hline
\end{tabular}
\end{center}

\subsection{Mit unbestimmtem Artikel}
\begin{center}
\begin{tabular}{|l|c|c|c|c|}
\hline
\textbf{Kasus} & \textbf{Maskulin} & \textbf{Feminin} & \textbf{Neutral} & \textbf{Plural} \\ \hline
Nominativ & ein guter & eine gute & ein gutes & keine guten \\ \hline
Akkusativ & einen guten & eine gute & ein gutes & keine guten \\ \hline
Dativ & einem guten & einer guten & einem guten & keinen guten \\ \hline
Genitiv & eines guten & einer guten & eines guten & keiner guten \\ \hline
\end{tabular}
\end{center}

\section{Relativpronomen (Anhang)}
\label{sec:RelativpronomenAnhang}

\begin{center}
\begin{tabular}{|l|c|c|c|c|}
\hline
\textbf{Kasus} & \textbf{Maskulin} & \textbf{Feminin} & \textbf{Neutral} & \textbf{Plural} \\ \hline
Nominativ & der & die & das & die \\ \hline
Akkusativ & den & die & das & die \\ \hline
Dativ & dem & der & dem & denen \\ \hline
Genitiv & dessen & deren & dessen & deren \\ \hline
\end{tabular}
\end{center}

\section{Fragepronomen}
\label{sec:Fragepronomen}

\begin{center}
\begin{tabular}{|l|c|c|c|c|}
\hline
\textbf{Kasus} & \textbf{Maskulin} & \textbf{Feminin} & \textbf{Neutral} & \textbf{Plural} \\ \hline
Nominativ & wer? & was? \\ \hline
Akkusativ & wen? & was? \\ \hline
Dativ & wem? \\ \hline
Genitiv & wessen? \\ \hline
\end{tabular}
\end{center}

\section{Wechselpräpositionen}
\label{sec:Wechselpraepositionen}

\begin{center}
\begin{tabular}{|l|c|c|}
\hline
\textbf{Präposition} & \textbf{Wo? (Dativ)} & \textbf{Wohin? (Akkusativ)} \\ \hline
in & in der Schule & in die Schule \\ \hline
an & an der Wand & an die Wand \\ \hline
auf & auf dem Tisch & auf den Tisch \\ \hline
unter & unter dem Tisch & unter den Tisch \\ \hline
vor & vor dem Haus & vor das Haus \\ \hline
hinter & hinter dem Haus & hinter das Haus \\ \hline
neben & neben dem Haus & neben das Haus \\ \hline
zwischen & zwischen den Bäumen & zwischen die Bäume \\ \hline
über & über dem Tisch & über den Tisch \\ \hline
\end{tabular}
\end{center}

\section{Übersicht: Präpositionen nach Kasus}
\label{sec:PraepositionenNachKasus}

\subsection{Dativ-Präpositionen}
\begin{center}
\begin{tabular}{|l|l|}
\hline
\textbf{Präposition} & \textbf{Bedeutung/Beispiel} \\ \hline
aus & Herkunft: Ich komme aus Berlin. \\ \hline
bei & Nähe: Ich bin bei Maria. \\ \hline
mit & Begleitung: Ich fahre mit dem Bus. \\ \hline
nach & Richtung: Nach dem Essen... \\ \hline
seit & Zeit: Seit zwei Jahren... \\ \hline
von & Herkunft: Das Buch ist von ihm. \\ \hline
zu & Richtung: Ich gehe zur Schule. \\ \hline
ab & Zeitpunkt: Ab morgen... \\ \hline
außer & Ausnahme: Außer mir war niemand da. \\ \hline
\end{tabular}
\end{center}

\subsection{Akkusativ-Präpositionen}
\begin{center}
\begin{tabular}{|l|l|}
\hline
\textbf{Präposition} & \textbf{Bedeutung/Beispiel} \\ \hline
bis & Richtung/Zeit: Bis morgen! \\ \hline
durch & Richtung: Durch den Wald \\ \hline
gegen & Richtung: Gegen die Wand \\ \hline
ohne & Ausschluss: Ohne dich \\ \hline
um & Richtung: Um die Ecke \\ \hline
\end{tabular}
\end{center}

\subsection{Genitiv-Präpositionen}
\begin{center}
\begin{tabular}{|l|l|}
\hline
\textbf{Präposition} & \textbf{Bedeutung/Beispiel} \\ \hline
trotz & trotz des Regens... \\ \hline
wegen & Wegen des Verkehrs... \\ \hline
während & Während des Films... \\ \hline
statt & Statt des Bruders... \\ \hline
außerhalb & Außerhalb der Stadt \\ \hline
innerhalb & Innerhalb der Woche \\ \hline
aufgrund & Aufgrund der Situation... \\ \hline
oberhalb & Oberhalb des Berges \\ \hline
unterhalb & Unterhalb des Tales \\ \hline
diesseits & Diesseits des Flusses \\ \hline
jenseits & Jenseits des Meeres \\ \hline
\end{tabular}
\end{center}

\section{Reflexive Verben mit Präpositionen (sich + Präposition)}
\label{sec:ReflexivPraep}

Viele reflexive Verben werden mit festen Präpositionen verwendet:

\subsection{sich an... gewöhnen}
\begin{center}
\begin{tabular}{|l|c|c|}
\hline
\textbf{Pronomen} & \textbf{Beispiel} & \textbf{Bedeutung} \\ \hline
ich & Ich gewöhne mich an das Wetter. & adapt to \\ \hline
du & Du gewöhnst dich an die neue Schule. & adapt to \\ \hline
er/sie/es & Er gewöhnt sich an den Job. & adapt to \\ \hline
wir & Wir gewöhnen uns an die Stadt. & adapt to \\ \hline
\end{tabular}
\end{center}

\subsection{sich auf... freuen}
\begin{center}
\begin{tabular}{|l|c|c|}
\hline
\textbf{Pronomen} & \textbf{Beispiel} & \textbf{Bedeutung} \\ \hline
ich & Ich freue mich auf den Urlaub. & look forward to \\ \hline
du & Du freust dich auf das Fest. & look forward to \\ \hline
er/sie/es & Sie freut sich auf dich. & look forward to \\ \hline
wir & Wir freuen uns auf Weihnachten. & look forward to \\ \hline
\end{tabular}
\end{center}

\subsection{sich auf... vorbereiten}
\begin{center}
\begin{tabular}{|l|c|c|}
\hline
\textbf{Pronomen} & \textbf{Beispiel} & \textbf{Bedeutung} \\ \hline
ich & Ich bereite mich auf die Prüfung vor. & prepare for \\ \hline
du & Du bereitest dich auf das Interview vor. & prepare for \\ \hline
er/sie/es & Er bereitet sich auf die Reise vor. & prepare for \\ \hline
wir & Wir bereiten uns auf die Präsentation vor. & prepare for \\ \hline
\end{tabular}
\end{center}

\subsection{sich über... ärgern}
\begin{center}
\begin{tabular}{|l|c|c|}
\hline
\textbf{Pronomen} & \textbf{Beispiel} & \textbf{Bedeutung} \\ \hline
ich & Ich ärgere mich über den Lärm. & be annoyed at \\ \hline
du & Du ärgerst dich über die Note. & be annoyed at \\ \hline
er/sie/es & Er ärgert sich über den Chef. & be annoyed at \\ \hline
wir & Wir ärgern uns über das Wetter. & be annoyed at \\ \hline
\end{tabular}
\end{center}

\subsection{sich über... informieren}
\begin{center}
\begin{tabular}{|l|c|c|}
\hline
\textbf{Pronomen} & \textbf{Beispiel} & \textbf{Bedeutung} \\ \hline
ich & Ich informiere mich über die Angebote. & inform about \\ \hline
du & Du informierst dich über das Programm. & inform about \\ \hline
er/sie/es & Er informiert sich über die Preise. & inform about \\ \hline
wir & Wir informieren uns über die Bedingungen. & inform about \\ \hline
\end{tabular}
\end{center}

\section{Nomen: Artikel und Plural}
\label{sec:NomenArtikel}

\subsection{Maskulin (der)}
\begin{center}
\begin{tabular}{|l|c|c|l|}
\hline
\textbf{Nomen} & \textbf{Artikel} & \textbf{Plural} & \textbf{Bedeutung} \\ \hline
der Mann & der & die Männer & man \\ \hline
der Junge & der & die Jungen & boy \\ \hline
der Vater & der & die Väter & father \\ \hline
der Sohn & der & die Söhne & son \\ \hline
der Bruder & der & die Brüder & brother \\ \hline
der Freund & der & die Freunde & friend \\ \hline
der Lehrer & der & die Lehrer & teacher \\ \hline
der Arzt & der & die Ärzte & doctor \\ \hline
der Student & der & die Studenten & student \\ \hline
der Mensch & der & die Menschen & human \\ \hline
der Herr & der & die Herren & gentleman \\ \hline
der Kunde & der & die Kunden & customer \\ \hline
der Kollege & der & die Kollegen & colleague \\ \hline
der Chef & der & die Chefs & boss \\ \hline
der Hund & der & die Hunde & dog \\ \hline
der Vogel & der & die Vögel & bird \\ \hline
der Tisch & der & die Tische & table \\ \hline
der Stuhl & der & die Stühle & chair \\ \hline
der Computer & der & die Computer & computer \\ \hline
der Tag & der & die Tage & day \\ \hline
\end{tabular}
\end{center}

\subsection{Feminin (die)}
\begin{center}
\begin{tabular}{|l|c|c|l|}
\hline
\textbf{Nomen} & \textbf{Artikel} & \textbf{Plural} & \textbf{Bedeutung} \\ \hline
die Frau & die & die Frauen & woman \\ \hline
die Mutter & die & die Mütter & mother \\ \hline
die Schwester & die & die Schwestern & sister \\ \hline
die Tochter & die & die Töchter & daughter \\ \hline
die Freundin & die & die Freundinnen & (female) friend \\ \hline
die Lehrerin & die & die Lehrerinnen & (female) teacher \\ \hline
die Ärztin & die & die Ärztinnen & (female) doctor \\ \hline
die Studentin & die & die Studentinnen & (female) student \\ \hline
die Katze & die & die Katzen & cat \\ \hline
die Blume & die & die Blumen & flower \\ \hline
die Schule & die & die Schulen & school \\ \hline
die Stadt & die & die Städte & city \\ \hline
die Uhr & die & die Uhren & clock/watch \\ \hline
die Lampe & die & die Lampen & lamp \\ \hline
die Tür & die & die Türen & door \\ \hline
die Hand & die & die Hände & hand \\ \hline
die Zeit & die & die Zeiten & time \\ \hline
die Arbeit & die & die Arbeiten & work \\ \hline
die Frage & die & die Fragen & question \\ \hline
die Sprache & die & die Sprachen & language \\ \hline
\end{tabular}
\end{center}

\subsection{Neutral (das)}
\begin{center}
\begin{tabular}{|l|c|c|l|}
\hline
\textbf{Nomen} & \textbf{Artikel} & \textbf{Plural} & \textbf{Bedeutung} \\ \hline
das Kind & das & die Kinder & child \\ \hline
das Buch & das & die Bücher & book \\ \hline
das Haus & das & die Häuser & house \\ \hline
das Auto & das & die Autos & car \\ \hline
das Bild & das & die Bilder & picture \\ \hline
das Wort & das & die Wörter & word \\ \hline
das Jahr & das & die Jahre & year \\ \hline
das Zimmer & das & die Zimmer & room \\ \hline
das Fenster & das & die Fenster & window \\ \hline
das Telefon & das & die Telefone & telephone \\ \hline
das Radio & das & die Radios & radio \\ \hline
das Foto & das & die Fotos & photo \\ \hline
das Baby & das & die Babys & baby \\ \hline
das Problem & das & die Probleme & problem \\ \hline
das Land & das & die Länder & country/land \\ \hline
das Meer & das & die Meere & sea \\ \hline
das Tier & das & die Tiere & animal \\ \hline
das Wasser & das & - & water \\ \hline
das Brot & das & - & bread \\ \hline
das Geld & das & - & money \\ \hline
\end{tabular}
\end{center}

\section{Nützliche Sätze für Meinungen und Optionen}
\label{sec:Meinungen}

\subsection{Meinungen ausdrücken}
\begin{center}
\begin{tabular}{|l|p{8cm}|}
\hline
\textbf{Satz} & \textbf{Bedeutung} \\ \hline
Ich finde, dass... & I think that... \\ \hline
Ich meine, dass... & I believe that... \\ \hline
Ich bin der Meinung, dass... & I am of the opinion that... \\ \hline
Meiner Meinung nach... & In my opinion... \\ \hline
Ich denke, dass... & I think that... \\ \hline
Ich halte es für... & I consider it to be... \\ \hline
Es scheint mir, dass... & It seems to me that... \\ \hline
Ich bin überzeugt, dass... & I am convinced that... \\ \hline
Es ist klar, dass... & It is clear that... \\ \hline
Natürlich... & Of course... \\ \hline
Klar... & Obviously... \\ \hline
Leider... & Unfortunately... \\ \hline
Ich stimme zu. & I agree. \\ \hline
Ich stimme nicht zu. & I disagree. \\ \hline
Das sehe ich anders. & I see it differently. \\ \hline
Das ist richtig. & That is correct. \\ \hline
Das ist falsch. & That is wrong. \\ \hline
\end{tabular}
\end{center}

\subsection{Optionen und Alternativen vorschlagen}
\begin{center}
\begin{tabular}{|l|p{8cm}|}
\hline
\textbf{Satz} & \textbf{Bedeutung} \\ \hline
Einerseits... andererseits... & On one hand... on the other hand... \\ \hline
Es gibt mehrere Möglichkeiten. & There are several possibilities. \\ \hline
Man könnte... & One could... \\ \hline
Man könnte auch... & One could also... \\ \hline
Wie wäre es mit...? & How about...? \\ \hline
Was hältst du von...? & What do you think of...? \\ \hline
Ich schlage vor, dass... & I suggest that... \\ \hline
Ich empfehle... & I recommend... \\ \hline
Es wäre besser, wenn... & It would be better if... \\ \hline
Anstelle von... könnte man... & Instead of... one could... \\ \hline
Alternativ könnte... & Alternatively... could... \\ \hline
Eine andere Möglichkeit wäre... & Another possibility would be... \\ \hline
Das beste wäre... & The best would be... \\ \hline
Im Gegensatz dazu... & In contrast to that... \\ \hline
\end{tabular}
\end{center}

\subsection{Zustimmen und Widersprechen}
\begin{center}
\begin{tabular}{|l|p{8cm}|}
\hline
\textbf{Satz} & \textbf{Bedeutung} \\ \hline
Ja, das stimmt. & Yes, that's right. \\ \hline
Genau so ist es. & Exactly. \\ \hline
Da hast du recht. & You're right there. \\ \hline
Ich bin ganz deiner Meinung. & I am completely of your opinion. \\ \hline
Das sehe ich genauso. & I see it the same way. \\ \hline
Nein, das sehe ich anders. & No, I see it differently. \\ \hline
Da muss ich widersprechen. & I have to disagree there. \\ \hline
Das ist nicht richtig. & That's not correct. \\ \hline
Ich sehe das nicht so. & I don't see it that way. \\ \hline
Das kann man nicht generalisieren. & You can't generalize that. \\ \hline
\end{tabular}
\end{center}

\section{Wichtige Wortverbindungen}
\label{sec:Wortverbindungen}

\subsection{Nomen + Präposition}
\begin{center}
\begin{tabular}{|l|c|l|}
\hline
\textbf{Nomen} & \textbf{+ Präposition} & \textbf{Beispiel} \\ \hline
die Angst & vor & Angst vor dem Hund \\ \hline
der Spaß & an & Spaß am Sport \\ \hline
das Interesse & an & Interesse an Musik \\ \hline
die Hoffnung & auf & Hoffnung auf Erfolg \\ \hline
die Freude & an & Freude am Lernen \\ \hline
der Ärger & über & Ärger über die Entscheidung \\ \hline
die Frage & nach & die Frage nach dem Sinn \\ \hline
der Grund & für & der Grund für die Verspätung \\ \hline
das Problem & mit & das Problem mit dem Computer \\ \hline
die Lösung & für & die Lösung für das Problem \\ \hline
\end{tabular}
\end{center}

\subsection{Adjektiv + Präposition}
\begin{center}
\begin{tabular}{|l|c|l|}
\hline
\textbf{Adjektiv} & \textbf{+ Präposition} & \textbf{Beispiel} \\ \hline
stolz & auf & stolz auf die Leistung \\ \hline
zufrieden & mit & zufrieden mit dem Ergebnis \\ \hline
glücklich & über & glücklich über die Nachricht \\ \hline
traurig & über & traurig über das Ende \\ \hline
wütend & auf & wütend auf den Chef \\ \hline
nervös & wegen & nervös wegen der Prüfung \\ \hline
abhängig & von & abhängig von den Eltern \\ \hline
interessiert & an & interessiert an Kunst \\ \hline
reich & an & reich an Erfahrungen \\ \hline
arm & an & arm an Zeit \\ \hline
\end{tabular}
\end{center}

\chapterend